\section{Feasibility Region}
\label{sec:cap3:feasibility-region}

After the optimization problem has been condensed into input variables and linear inequalities, it is useful to characterize the states for which the problem is feasible. For a given current state \highlight{$x[k]$}, feasibility means that there exists at least one predicted input sequence \highlight{$\mathbf{\hat{u}}[k]$} satisfying all input and state constraints over the prediction horizon.

Using the condensed constraint form obtained in \eqref{eq:cap3:optimization-constraints}, the feasibility region is defined as

\noindent
\begin{equation}
    \mathcal{D} =
    \left\{
        x[k] \in \mathcal{X} :
        \exists \mathbf{\hat{u}}[k] \in \mathbb{R}^{pN}
        \text{ such that }
        \mathbf{S}\mathbf{\hat{u}}[k] \le \mathbf{b}(x[k])
    \right\}.
\end{equation}

\noindent
Thus, if \highlight{$x[k] \in \mathcal{D}$}, the finite-horizon MPC problem admits at least one feasible input sequence. If \highlight{$x[k] \notin \mathcal{D}$}, the constraints cannot be satisfied for the selected prediction horizon.

The same set can be computed by introducing the joint vector

\noindent
\begin{equation}
    w =
    \left[
        \begin{matrix}
            \mathbf{\hat{u}}[k] \\
            x[k]
        \end{matrix}
    \right].
\end{equation}

\noindent
The condensed state and input constraints can then be written as a polyhedron in the joint space of input sequences and current states:

\noindent
\begin{equation}
    S_w w \le b_w,
\end{equation}

\noindent
with

\noindent
\begin{equation}
    S_w =
    \left[
        \begin{matrix}
            \mathbf{S_x} \mathbf{B} & \mathbf{S_x} \mathbf{A} \\
            \mathbf{S_u} & 0 \\
            0 & S_x
        \end{matrix}
    \right],
    \quad
    b_w =
    \left[
        \begin{matrix}
            \mathbf{b_x} \\
            \mathbf{b_u} \\
            b_x
        \end{matrix}
    \right].
\end{equation}

\noindent
The last block row enforces \highlight{$x[k] \in \mathcal{X}$}. The first two block rows enforce the predicted state and input constraints over the horizon.

The polyhedron \highlight{$\mathcal{P}_w = \{w : S_w w \le b_w\}$} contains all pairs \highlight{$(\mathbf{\hat{u}}[k], x[k])$} that satisfy the constraints. The feasibility region \highlight{$\mathcal{D}$} is obtained by projecting \highlight{$\mathcal{P}_w$} onto the state coordinates, which eliminates \highlight{$\mathbf{\hat{u}}[k]$} and keeps only the states for which at least one admissible input sequence exists.

The projection can be performed using the \textit{projection} function from the \textit{MPT Toolbox}.
